%% Correction audit de veracite - 2026-09-20 : citations de rapport de jury sans
%% support textuel dans le corpus (rapports 2018-2025) remplacees par des verbatims
%% verifies ou supprimees. Detail ci-dessous.
%% l.306 (ex) : citation "jury 2024" reformulee (lois fondamentales / composants
%%   elementaires) INTROUVABLE telle quelle -> remplacee par RJ2024 l.511 (verbatim).
%% l.307 (ex) : citation "jury 2025" reformulee et completee d'une phrase ("Dessiner les
%%   schemas equivalents pour chaque intervalle de commutation") INTROUVABLE
%%   -> remplacee par le verbatim exact de RJ2025 l.486.
\expandafter\ifx\csname CAPETINCLUDEMODE\endcsname\relax
\documentclass[11pt,a4paper]{article}
% preamble.tex — packages et commandes communes a toutes les fiches CAPET SII IE
\input{glyphtounicode}
\pdfgentounicode=1
\usepackage{cmap}
\usepackage[T1]{fontenc}
\usepackage[utf8]{inputenc}
\usepackage[french]{babel}
% babel[french] active a \begin{document} des raccourcis de ponctuation (:;!?)
% qui cassent la syntaxe to[...] de circuitikz (erreur "+ or - expected") ;
% on les desactive, ce qui ne change que l'espacement fin avant ces signes.
\AtBeginDocument{\shorthandoff{:;!?}}
\usepackage[a4paper, top=1.6cm, bottom=1.6cm, left=1.4cm, right=1.4cm]{geometry}
\setlength{\headheight}{14.5pt}
\usepackage{amsmath, amssymb, amsfonts}
\usepackage{siunitx}
\sisetup{locale=FR, output-decimal-marker={,}, inter-unit-product=\ensuremath{{}\cdot{}},
         group-minimum-digits=4}
% Unites non fournies par siunitx, declarees ici une fois pour toutes (fiches + exercices).
% Les gardes \ifdefined evitent tout conflit avec une redeclaration locale restee dans un
% fichier : ne pas retirer les gardes, ne pas redeclarer ces unites dans les fichiers.
\DeclareSIUnit{\jour}{j}
\ifdefined\var\else\DeclareSIUnit{\var}{var}\fi   % puissance reactive
\ifdefined\vad\else\DeclareSIUnit{\vad}{VAD}\fi   % puissance deformante
\ifdefined\tr\else\DeclareSIUnit{\tr}{tr}\fi      % tour (tr/min)
\usepackage{array, booktabs, multirow, tabularx}
\usepackage{xcolor}
\usepackage[most, breakable]{tcolorbox}
\usepackage{multicol}
\usepackage{enumitem}
\usepackage{mathtools}
\usepackage{lmodern}
% Refonte visuelle 2026 (direction "moderne coloré") : Lato pour le texte,
% maths sans-serif assorties via newtxsf (charge apres amsmath/amssymb).
\usepackage[default]{lato}
\usepackage{newtxsf}
\usepackage{microtype}
\usepackage{fancyhdr}
\usepackage{tikz}
\usetikzlibrary{positioning,arrows.meta}
\usepackage{circuitikz}
\usepackage{lastpage}

% ── Couleurs par domaine (palette modernisée 2026) ────────────────────
\definecolor{colElec}{RGB}{29, 78, 216}        % Électrotechnique — bleu vif
\definecolor{colElecLight}{RGB}{219, 234, 254} % fond bleu clair
\definecolor{colAuto}{RGB}{21, 128, 61}        % Automatique — vert
\definecolor{colMeca}{RGB}{194, 65, 12}        % Mécanique — orange
\definecolor{colTherm}{RGB}{185, 28, 28}       % Thermique — rouge
\definecolor{colQT}{RGB}{232, 244, 233}        % (hérité, plus utilisé par questiontype)
\definecolor{colJury}{RGB}{255, 248, 225}      % fond jury (crème)
\definecolor{colJuryFrame}{RGB}{202, 138, 4}   % bordure/bandeau jury (ambre)
\definecolor{colGrayLight}{RGB}{245, 245, 247} % fond gris neutre
\definecolor{colSeparator}{RGB}{180, 180, 180} % gris séparateur
\definecolor{colExo}{RGB}{8, 145, 178}          % Banque d'exercices — cyan (distinct des fiches)
% ── Encre et gris de la refonte ───────────────────────────────────────
\definecolor{colInk}{RGB}{23, 32, 42}          % texte courant (quasi-noir)
\definecolor{colSlate}{RGB}{83, 98, 118}       % gris ardoise (méta, en-têtes)
\definecolor{colRule}{RGB}{203, 213, 225}      % filets clairs
\definecolor{colJuryBg}{RGB}{255, 248, 225}    % fond crème jury
\definecolor{colPresence}{RGB}{83, 98, 118}    % bloc 5 (présence dans les sujets)
\AtBeginDocument{\color{colInk}}

% ── Style encadré exercice + corrigé (banque d'exercices) ─────────────
\tcbset{
  exoBox/.style={
    breakable,
    enhanced,
    colback=colExo!8!white,
    colframe=colExo,
    colbacktitle=colExo,
    coltitle=white,
    fonttitle=\bfseries\small,
    left=4pt, right=4pt, top=3pt, bottom=3pt,
    before skip=5pt, after skip=5pt,
  }
}

% ── Bascule enonce / corrige (banque d'exercices) ─────────────────────
% \CAPETCORRIGE est defini en ligne de commande (pdflatex "\def\CAPETCORRIGE{1}\input{...}")
% pour compiler la variante corriges d'un meme fichier source ; par defaut (non defini),
% la compilation produit la variante enonces. \CAPETSHOW{enonce}{corrige} choisit la branche.
%
% TROIS MODES (refonte 2026-09-20, manuel de revision) :
%   - aucun drapeau   -> imprime l'enonce            (variante enonces, INCHANGE)
%   - \CAPETCORRIGE   -> imprime le corrige          (variante corriges, INCHANGE)
%   - \CAPETMANUEL    -> imprime l'enonce ET met le corrige de cote pour l'annexe
%
% Le mode manuel stocke chaque corrige dans une macro (\csgdef), sans ecriture
% fichier ni re-tokenisation : le contenu (circuitikz, accents, %) traverse intact.
% L'annexe les rejoue ensuite dans l'ordre via \CAPETANNEXE.
\usepackage{etoolbox}

\newcounter{CAPETexo}          % numero continu d'exercice dans le manuel
\setcounter{CAPETexo}{0}
\newcounter{CAPETannexe}       % compteur de relecture en annexe

% Titre courant de la notion en cours (pose par \CAPETNOTION), sert d'intertitre annexe
\newcommand{\CAPETnotioncode}{}
\newcommand{\CAPETnotiontitre}{}

% \CAPETNOTION{code}{titre} : ouvre une notion dans le manuel (memorise le libelle,
% et marque le rang du premier exercice de la notion pour l'intertitre d'annexe).
\newcommand{\CAPETNOTION}[2]{%
  \renewcommand{\CAPETnotioncode}{#1}%
  \renewcommand{\CAPETnotiontitre}{#2}%
  \csgdef{CAPETnotionstart\the\numexpr\value{CAPETexo}+1\relax}{#1 -- #2}%
}

% ── \degree robuste ───────────────────────────────────────────────────
% Certains fichiers d'exercices utilisent \degree. Ni le preambule ni siunitx
% ne le definissent, et \providecommand local ne suffit pas : sous hyperref,
% \degree existe deja comme commande factice (\ERROR), donc \providecommand
% s'abstient et l'appel echoue a la compilation du manuel (erreur invisible
% en compilation isolee). On le (re)definit donc inconditionnellement ici,
% apres le chargement des packages, pour tout le corpus.
\AtBeginDocument{\providecommand{\degree}{}\renewcommand{\degree}{\ensuremath{^\circ}}}

% Libelle decoratif : les fichiers d'exercices utilisent aussi \CAPETSHOW dans
% leur \fancyhead et leur bloc de titre, pour afficher "(enonces)" ou "(corriges)".
% Ces appels-la ne sont PAS des exercices : en mode manuel ils ne doivent ni
% consommer un numero, ni alimenter l'annexe. On les reconnait a leur premier
% argument, strictement egal au libelle ci-dessous.
\newcommand{\CAPETLIBELLEENONCES}{(\'enonc\'es)}

\ifdefined\CAPETMANUEL
  % --- Mode manuel : enonce en place, corrige differe vers l'annexe ---
  \newcommand{\CAPETSHOW}[2]{%
    \def\CAPETargun{#1}%
    \ifx\CAPETargun\CAPETLIBELLEENONCES
      #1%            appel decoratif : rendu tel quel, hors numerotation
    \else
      \CAPETSHOWEXO{#1}{#2}%
    \fi
  }
  \newcommand{\CAPETSHOWEXO}[2]{%
    \stepcounter{CAPETexo}%
    \csgdef{CAPETcorrige\theCAPETexo}{#2}%
    \edef\CAPETtmpnum{\theCAPETexo}%
    #1%
    \par\vspace{2pt}%
    {\footnotesize\itshape\color{colExo}%
     \hyperref[capet:cor:\CAPETtmpnum]{Corrige de l'exercice \CAPETtmpnum\ $\rightarrow$ p.~\pageref{capet:cor:\CAPETtmpnum}}%
     \label{capet:eno:\CAPETtmpnum}}%
  }
  % \CAPETANNEXE : rejoue tous les corriges captures, dans l'ordre, avec intertitres
  \newcommand{\CAPETANNEXE}{%
    \setcounter{CAPETannexe}{0}%
    \loop\ifnum\value{CAPETannexe}<\value{CAPETexo}%
      \stepcounter{CAPETannexe}%
      \edef\CAPETcurnum{\theCAPETannexe}%
      \ifcsname CAPETnotionstart\CAPETcurnum\endcsname
        \par\medskip
        {\normalsize\bfseries\color{colExo}\csuse{CAPETnotionstart\CAPETcurnum}}%
        \par\nobreak\smallskip
      \fi
      \begin{tcolorbox}[exoBox, breakable,
        title={Corrige de l'exercice \CAPETcurnum\
               \hfill\normalfont\footnotesize
               \hyperref[capet:eno:\CAPETcurnum]{enonce p.~\pageref{capet:eno:\CAPETcurnum}}}]%
      \label{capet:cor:\CAPETcurnum}%
      \csuse{CAPETcorrige\CAPETcurnum}%
      \end{tcolorbox}%
    \repeat
  }
\else
  % --- Modes historiques, comportement strictement inchange ---
  \newcommand{\CAPETSHOW}[2]{\ifdefined\CAPETCORRIGE #2\else #1\fi}
  \newcommand{\CAPETANNEXE}{}
\fi

% ── Difficulté 1 à 3 (étoiles pleines, amssymb) ───────────────────────
\newcommand{\difficulte}[1]{%
  \ifcase#1\relax
  \or $\bigstar$\;niveau 1/3%
  \or $\bigstar\bigstar$\;niveau 2/3%
  \or $\bigstar\bigstar\bigstar$\;niveau 3/3%
  \fi
}

% ── Style encadré extrait de document technique (banque d'exercices) ──
\tcbset{
  docTech/.style={
    breakable,
    enhanced,
    colback=colGrayLight,
    colframe=colSeparator,
    fonttitle=\bfseries\small\itshape,
    coltitle=black,
    colbacktitle=colGrayLight,
    left=4pt, right=4pt, top=3pt, bottom=3pt,
    before skip=4pt, after skip=4pt,
  }
}

% ── Style tcolorbox de section principale ────────────────────────────
% Refonte 2026 : bandeau titre coloré arrondi + fond très pâle, sans cadre.
\tcbset{
  ficheSection/.style={
    breakable, enhanced jigsaw,
    frame hidden, boxrule=0pt, arc=1.8mm,
    colback=#1!4!white,
    colbacktitle=#1,
    coltitle=white,
    fonttitle=\sffamily\bfseries,
    left=7pt, right=7pt, top=5pt, bottom=5pt,
    toptitle=3pt, bottomtitle=3pt,
    before skip=10pt, after skip=8pt,
  }
}

% ── Style encadré question type ──────────────────────────────────────
% Carte blanche, bandeau sombre, liseré très clair.
\tcbset{
  questiontype/.style={
    breakable, enhanced jigsaw,
    colback=white,
    frame hidden, boxrule=0pt, arc=1.8mm,
    colbacktitle=colElec!78!black,
    coltitle=white,
    fonttitle=\sffamily\bfseries\small,
    left=7pt, right=7pt, top=4pt, bottom=4pt,
    toptitle=2.5pt, bottomtitle=2.5pt,
    before skip=8pt, after skip=8pt,
    borderline={0.6pt}{0pt}{colElec!12!white},
  }
}

% ── Style encadré jury ───────────────────────────────────────────────
% Bandeau ambre, fond crème, sans cadre.
\tcbset{
  jury/.style={
    breakable, enhanced jigsaw,
    colback=colJuryBg,
    frame hidden, boxrule=0pt, arc=1.8mm,
    colbacktitle=colJuryFrame,
    coltitle=white,
    fonttitle=\sffamily\bfseries,
    left=7pt, right=7pt, top=4pt, bottom=4pt,
    toptitle=3pt, bottomtitle=3pt,
    before skip=10pt, after skip=8pt,
  }
}

% ── En-tête / pied de page communs (refonte 2026) ─────────────────────
% #1 code fiche, #2 titre court, #3 couleur domaine
\newcommand{\ficheHeaders}[3]{%
  \fancyhead[L]{\small\sffamily{\bfseries\textcolor{#3}{#1}}\hspace{0.7em}\textcolor{colSlate}{#2}}%
  \fancyhead[R]{\small\sffamily\textcolor{colSlate}{CAPET SII IE\, --- \,Session 2026}}%
  \fancyfoot[C]{\small\sffamily\textcolor{colSlate}{\thepage\,/\,\pageref{LastPage}}}%
  \fancyhead[C]{}\fancyfoot[L]{}\fancyfoot[R]{}%
}

% ── Bloc en-tête de fiche (BLOC 0, refonte 2026) ──────────────────────
% #1 code  #2 titre  #3 couleur domaine  #4 domaine  #5 ligne présence  #6 sous-thèmes
\newcommand{\ficheEntete}[6]{%
  \begin{tcolorbox}[enhanced, breakable,
    frame hidden, boxrule=0pt, arc=2.5mm,
    interior style={left color=#3, right color=#3!75!black},
    left=9pt, right=9pt, top=7pt, bottom=7pt,
    before skip=0pt, after skip=6pt]
    {\color{white}\sffamily\bfseries\LARGE #1\quad #2}\\[4pt]
    {\color{#3!22!white}\sffamily\small \MakeUppercase{#4}\quad\textbar\quad #5}\\[1pt]
    {\color{#3!35!white}\sffamily\footnotesize Sous-thèmes : #6}
  \end{tcolorbox}%
}

% ── Mise en page ─────────────────────────────────────────────────────
\pagestyle{fancy}
\fancyhf{}
\renewcommand{\headrulewidth}{0.5pt}
\renewcommand{\footrulewidth}{0pt}
\AtBeginDocument{%
  \renewcommand{\headrule}{{\color{colRule}\hrule width\headwidth height 0.5pt}}%
}

\setlist[itemize]{leftmargin=1.2em, itemsep=1pt, topsep=2pt, parsep=0pt}
\setlist[enumerate]{leftmargin=1.2em, itemsep=3pt, topsep=2pt, parsep=0pt}

\setlength{\parskip}{2pt}
\setlength{\columnsep}{10pt}

\def\CAPETBEGINFILE{\begin{document}}
\def\CAPETENDFILE{\end{document}}
\else
\def\CAPETBEGINFILE{}
\def\CAPETENDFILE{}
\fi

\ficheHeaders{T18}{Composants électroniques}{colElec}

\CAPETBEGINFILE

% -----------------------------------------------------------------------
%  BLOC EN-TETE
% -----------------------------------------------------------------------

\ficheEntete{T18}{Circuits électroniques - composants (diode, transistor)}{colElec}%
  {électronique - Composants et convertisseurs}%
  {\textbf{Présence :} \textbf{4/13} sujets (2013-2026) - thème \emph{occasionnel}}%
  {diode (modèle idéal, seuil, droite de charge) - transistor BJT (3 régimes) - MOSFET et IGBT - hacheur abaisseur (courant inductance, ondulation) - diode de roue libre - rapport cyclique}

% -----------------------------------------------------------------------
%  1. RAPPELS DE COURS
% -----------------------------------------------------------------------

\begin{tcolorbox}[
  ficheSection=colElec,
  title={1. Rappels de cours},
]

% -- 1.1 Diode --------------------------------------------------------
\textbf{1.1 La diode}

\smallskip
\begin{center}
\renewcommand{\arraystretch}{1.4}
\begin{tabular}{lll}
\toprule
\textbf{Modèle} & \textbf{État conducteur ($V_{AK}>0$)} & \textbf{État bloqué ($V_{AK}<0$)} \\
\midrule
Idéal & $V_{AK} = 0$, $i \geq 0$ & $i = 0$ \\
Avec seuil & $V_{AK} = V_{seuil}$ (Si : 0{,}6\,V) & $i = 0$ \\
\bottomrule
\end{tabular}
\end{center}

\smallskip
\textbf{Droite de charge et point de fonctionnement Q :}
\begin{center}
\begin{tikzpicture}[scale=0.9, font=\footnotesize]
  \draw[->] (0,0) -- (5.2,0) node[right] {$V_{AK}$};
  \draw[->] (0,-0.3) -- (0,3.2) node[above] {$i$};
  \draw[thick, colElec] (0,2.8) -- (4.6,0);
  \node[left] at (0,2.8) {\scriptsize $E/R$};
  \node[below] at (4.6,0) {\scriptsize $E$};
  \draw[thick, colElec!55, domain=0.9:1.3, smooth, samples=30] plot (\x, {(\x-0.9)*8});
  \draw[thick, colElec!55] (1.3,2.9) -- (1.3,3.15);
  \fill (1.05,1.2) circle (1.5pt);
  \node[right] at (1.12,1.2) {\scriptsize $Q$};
\end{tikzpicture}
\end{center}
\textbf{Droite de charge} pour un circuit $E = R\,i + V_{AK}$ : pente $= -1/R$, intersection axe $i$ en $E/R$ et axe $V$ en $E$. Le \emph{point de fonctionnement} Q est l'intersection avec la caractéristique $i(V_{AK})$.

\textbf{Diode de roue libre} : placée en anti-parallèle avec une charge inductive. À l'ouverture de l'interrupteur, elle assure la continuité du courant dans l'inductance et protège le composant actif contre les surtensions $L\,di/dt$.

\medskip\hrule\vspace{4pt}

% -- 1.2 Transistor bipolaire (BJT) -----------------------------------
\textbf{1.2 Transistor bipolaire (BJT) - trois régimes}

\smallskip
\begin{center}
\renewcommand{\arraystretch}{1.3}
\begin{tabular}{llll}
\toprule
\textbf{Régime} & \textbf{Condition} & \textbf{Courant $I_C$} & \textbf{Usage} \\
\midrule
Bloqué & $V_{BE} < 0{,}6\,\mathrm{V}$ & $\approx 0$ & Interrupteur ouvert \\
Actif (linéaire) & $V_{BE} \approx 0{,}6\,\mathrm{V}$, $V_{CE} > V_{CE(sat)}$ & $\beta I_B$ & Amplification \\
Saturé & $V_{BE} \approx 0{,}6\,\mathrm{V}$, $V_{CE} \approx 0{,}2\,\mathrm{V}$ & Limité par circuit & Interrupteur fermé \\
\bottomrule
\end{tabular}
\end{center}
En saturation, $I_C = \beta I_B$ n'est \emph{plus} valable.

\medskip\hrule\vspace{4pt}

% -- 1.3 MOSFET, IGBT et hacheur abaisseur ----------------------------
\textbf{1.3 MOSFET et IGBT - hacheur abaisseur}

\smallskip
\begin{center}
\renewcommand{\arraystretch}{1.3}
\begin{tabular}{lll}
\toprule
\textbf{Composant} & \textbf{Commande} & \textbf{Usage} \\
\midrule
MOSFET & Tension $V_{GS}$ & Haute fréquence, BT \\
IGBT & Tension $V_{GE}$ (+ diode anti-parallèle) & Variateurs, onduleurs HT \\
\bottomrule
\end{tabular}
\end{center}

\textbf{Hacheur abaisseur (Buck) - états :}
\begin{itemize}[itemsep=1pt]
  \item $t \in [0\,;\,\alpha T]$ : $Q_1$ fermé, $D_1$ bloquée - $v_{Q1} = 0$, $v_{D1} = -V_{DC}$, $v_L = V_{DC} - V_s$ - $i_L$ croît
  \item $t \in [\alpha T\,;\,T]$ : $Q_1$ ouvert, $D_1$ passante (roue libre) - $v_{Q1} = V_{DC}$, $v_{D1} = 0$, $v_L = -V_s$ - $i_L$ décroît
\end{itemize}
\[
  \langle V_{sortie} \rangle = \alpha \cdot V_{DC}
  \qquad;\qquad
  \Delta i_L = \frac{V_{DC}\,\alpha\,(1-\alpha)\,T}{L}
  \qquad;\qquad
  \Delta i_{L,max} \text{ pour } \alpha = 0{,}5
\]

\end{tcolorbox}

% -----------------------------------------------------------------------
%  2. FORMULAIRE SYNTHETIQUE
% -----------------------------------------------------------------------

\begin{tcolorbox}[
  ficheSection=colElec,
  title={2. Formulaire synthétique},
]
\small
\begin{multicols}{2}
\renewcommand{\arraystretch}{1.45}
\begin{tabular}{@{}p{4.5cm}ll@{}}
\toprule
\textbf{Formule} & \textbf{Usage} & \textbf{U.} \\
\midrule
\multicolumn{3}{@{}l}{\footnotesize\textit{Diode}} \\
$V_{seuil,Si} = 0{,}6\,\mathrm{V}$ & Seuil & V \\
$E = Ri + V_{AK}$ & Droite de charge & - \\
\midrule
\multicolumn{3}{@{}l}{\footnotesize\textit{BJT}} \\
$I_C = \beta I_B$ (régime actif) & Amplification & A \\
$V_{CE,sat} \approx 0{,}2\,\mathrm{V}$ & Saturation & V \\
\midrule
\multicolumn{3}{@{}l}{\footnotesize\textit{Hacheur}} \\
$\langle V_s \rangle = \alpha V_{DC}$ & Tension moy. & V \\
$\Delta i_L = V_{DC}\alpha(1-\alpha)T/L$ & Ondulation & A \\
\bottomrule
\end{tabular}

\columnbreak

\renewcommand{\arraystretch}{1.45}
\begin{tabular}{@{}p{4.5cm}ll@{}}
\toprule
\textbf{Formule} & \textbf{Usage} & \textbf{U.} \\
\midrule
\multicolumn{3}{@{}l}{\footnotesize\textit{Courant inductance}} \\
$i_L(t) = i_0 + (V_{DC}-V_s)t/L$ & $Q_1$ fermé & A \\
$i_L(t) = i_{max} - V_s(t-\alpha T)/L$ & $D_1$ passante & A \\
\midrule
\multicolumn{3}{@{}l}{\footnotesize\textit{Inductance min.}} \\
$L_{min} = V_{DC}\alpha(1-\alpha)T/\Delta i_{max}$ & Contin. condu. & H \\
\midrule
\multicolumn{3}{@{}l}{\footnotesize\textit{Données typiques}} \\
$V_{DC} = 750\,\mathrm{V}$ (traction 2019) & Bus DC & V \\
$V_{DC} = 48\,\mathrm{V}$ (Bakus 2023) & Bus DC & V \\
$\beta \approx 50\text{ à }500$ & BJT & - \\
\bottomrule
\end{tabular}
\end{multicols}
\normalsize
\end{tcolorbox}

% -----------------------------------------------------------------------
%  3. QUESTIONS TYPES
% -----------------------------------------------------------------------

\begin{tcolorbox}[
  ficheSection=colElec,
  title={3. Questions types - extraites des sujets réels},
]

% -- QT1 -------------------------------------------------------------

\begin{tcolorbox}[questiontype,
  title={QT1 - Courant dans l'inductance d'un hacheur abaisseur \hfill\normalfont\footnotesize E2 disc. 2025 (Q27-Q28)}
]
\textbf{Énoncé.}
Hacheur abaisseur avec $Q_1$ (IGBT) et $D_1$ (roue libre). $V_{DC} = 48\,\mathrm{V}$, $V_s = 24\,\mathrm{V}$, $L = 2\,\mathrm{mH}$, $\alpha = 0{,}6$, $T = 100\,\mu\mathrm{s}$.
Exprimer $i_{L1}(t)$ sur chaque intervalle par les lois de Kirchhoff.

\tcbline

\textbf{Corrigé.}

\textbf{Étape 1 — Intervalle $t \in [0\,;\,\alpha T]$ ($Q_1$ fermé, $D_1$ bloquée).}
Maille : $V_{DC} = v_L + V_s \Rightarrow v_L = V_{DC} - V_s = 48 - 24 = 24\,\mathrm{V}$
\[
  v_L = L\,\frac{di_L}{dt} \;\Rightarrow\; i_L(t) = i_0 + \frac{V_{DC} - V_s}{L}\,t = i_0 + \frac{24}{2\times10^{-3}}\,t \quad\text{(croissant)}
\]

\textbf{Étape 2 — Intervalle $t \in [\alpha T\,;\,T]$ ($Q_1$ ouvert, $D_1$ passante).}
Maille : $0 = v_L + V_s \Rightarrow v_L = -V_s = -24\,\mathrm{V}$
\[
  i_L(t) = i_{max} - \frac{V_s}{L}\,(t - \alpha T) = i_{max} - \frac{24}{2\times10^{-3}}\,(t - 60\,\mu\mathrm{s}) \quad\text{(décroissant)}
\]
Ondulation : $\Delta i_L = \frac{48 \times 0{,}6 \times (1-0{,}6) \times 100\times10^{-6}}{2\times10^{-3}} = \mathbf{0{,}576\,\mathrm{A}}$

\textbf{Conclusion.} Les expressions temporelles et l'ondulation calculée sont cohérentes avec un fonctionnement en conduction continue : le dimensionnement de $L$ est conforme au cahier des charges de commande du hacheur.
\end{tcolorbox}

\vspace{4pt}

% -- QT2 -------------------------------------------------------------

\begin{tcolorbox}[questiontype,
  title={QT2 - Rapport cyclique d'ondulation maximale et inductance minimale \hfill\normalfont\footnotesize E2 disc. 2023 (Q C6, robot Bakus)}
]
\textbf{Énoncé.}
Hacheur abaisseur : $V_{DC} = 48\,\mathrm{V}$, $T = 100\,\mu\mathrm{s}$, ondulation maximale admissible $\Delta i_{max} = 4\,\mathrm{A}$.
(a) Trouver le rapport cyclique $\alpha$ qui maximise $\Delta i_L$. (b) Calculer l'inductance minimale $L_{min}$.

\tcbline

\textbf{Corrigé.}

\textbf{Étape 1 (a) — Rapport cyclique critique.} $\Delta i_L = V_{DC}\,\alpha(1-\alpha)T/L$. Maximiser $\alpha(1-\alpha)$ - dériver par rapport à $\alpha$ et annuler :
\[
  \frac{d}{d\alpha}[\alpha(1-\alpha)] = 1 - 2\alpha = 0 \;\Rightarrow\; \alpha^* = \mathbf{0{,}5}
\]

\textbf{Étape 2 (b) — Inductance minimale.} En $\alpha = 0{,}5$ :
\[
  L_{min} = \frac{V_{DC} \times 0{,}5 \times 0{,}5 \times T}{\Delta i_{max}} = \frac{48 \times 0{,}25 \times 100 \times 10^{-6}}{4} = \frac{1{,}2 \times 10^{-3}}{4} = \mathbf{0{,}3\,\mathrm{mH}}
\]

\textbf{Conclusion.} $L_{min} = 0{,}3\,\mathrm{mH}$ n'est valable que pour le pire cas $\alpha = 0{,}5$ : en pratique on choisit $L$ strictement supérieure à cette valeur pour garantir la conduction continue quel que soit le point de fonctionnement réel.
\end{tcolorbox}

\vspace{4pt}

% -- QT3 -------------------------------------------------------------

\begin{tcolorbox}[questiontype,
  title={QT3 - Diode de protection - entrée analogique d'un CAN \hfill\normalfont\footnotesize E2 disciplinaire 2022 (Q26-Q27)}
]
\textbf{Énoncé.}
La charge des accumulateurs est mesurée au travers d'une résistance de shunt $R_0$, dont la tension image est appliquée sur l'entrée analogique ADC1 du CAN ($V_{ref} = \SI{3,5}{\volt}$). Une diode $D_1$ (Si) protège cette entrée. Expliquer le rôle de $R_0$, puis le principe de la protection.

\tcbline

\textbf{Corrigé.}

\textbf{Étape 1 — Rôle de $R_0$.} $R_0$ est une \textbf{résistance de shunt} : le courant des accumulateurs la traverse et y développe une tension \emph{image de ce courant}, seule grandeur exploitable par une entrée du CAN, qui ne mesure que des tensions.

\textbf{Étape 2 — Charge des accumulateurs.} Le courant circule de la borne $+$ vers la borne $-$~: la tension aux bornes du shunt est \textbf{positive}, la diode est \textbf{bloquée} et n'a aucune influence sur la mesure.

\textbf{Étape 3 — Décharge : le cas dangereux.} Lorsque les accumulateurs débitent vers le système, le courant \textbf{change de sens} : la tension appliquée sur l'entrée devient \textbf{négative}, ce qui est nuisible pour le CAN (mise en conduction des diodes de structure du microcontrôleur). La diode devient alors passante ($V_{AK} > \SI{0,6}{\volt}$) et écrête cette tension :
\[
  V_{ADC1} \geq -V_{seuil} = \SI{-0,6}{\volt}
\]

\textbf{Conclusion.} La protection joue donc \emph{par le bas}, contre une tension \textbf{négative} née de l'inversion du courant, et non contre une surtension : l'entrée du CAN ne peut pas descendre sous $\SI{-0,6}{\volt}$, conformément à l'exigence de robustesse du cahier des charges.
\end{tcolorbox}

\end{tcolorbox}

% -----------------------------------------------------------------------
%  4. ATTENTES DU JURY
% -----------------------------------------------------------------------

\begin{tcolorbox}[
  jury,
  title={4. Attentes du jury - Erreurs récurrentes et points de vigilance},
]

\begin{enumerate}[label=\textbf{\arabic*.}, leftmargin=*, itemsep=5pt]

  \item \textbf{Oublier la diode de roue libre.}
  Sans $D_1$ en anti-parallèle avec la charge inductive, la tension aux bornes du transistor diverge à l'ouverture ($v = L\,di/dt \to \infty$) et le détruit. Toujours justifier sa présence.

  \item \textbf{Ondulation maximale à $\alpha = 0{,}5$.}
  $\Delta i_L \propto \alpha(1-\alpha)$ est maximisée pour $\alpha = 0{,}5$, pas pour $\alpha = 1$. Erreur classique dans la formule.

  \item \textbf{Rapport cyclique et tension moyenne.}
  $\langle V_s \rangle = \alpha V_{DC}$ n'est valable qu'en \emph{conduction continue} ($i_L > 0$ en permanence). En conduction discontinue, la relation change.

  \item \textbf{Saturation du BJT $\neq$ régime actif.}
  En saturation, $I_C = \beta I_B$ n'est plus valable et $V_{CE} \approx 0{,}2\,\mathrm{V}$. Le transistor se comporte comme un interrupteur fermé.

  \item \textbf{Signe de $V_{AK}$ pour la diode.}
  En passant, $V_{AK} > 0$. En bloquée, $V_{AK} < 0$. L'erreur de signe est la faute la plus fréquente.

  \item \textbf{Schémas équivalents pour chaque intervalle.}
  Le jury attend que le candidat dessine le circuit équivalent ($Q_1$ ouvert ou fermé, $D_1$ passante ou bloquée) avant d'appliquer les lois de Kirchhoff.

\end{enumerate}

\medskip
\begin{tcolorbox}[
  colback=colJuryFrame!12!colJury,
  colframe=colJuryFrame,
  left=4pt, right=4pt, top=2pt, bottom=2pt,
  fontupper=\small\itshape
]
Rapport jury 2024 : «~Le jury regrette le manque de maîtrise des formules de base sur les réseaux triphasés, le régime sinusoïdal et les relations d'électrocinétique telles que la loi des mailles, la loi des nœuds ou encore la loi d'Ohm.~»\\
Rapport jury 2025 : «~Des résultats donnés directement, sans calcul et sans justification ne répondent pas aux attendus de l'épreuve.~»
\end{tcolorbox}

\end{tcolorbox}

% -----------------------------------------------------------------------
%  PIED - PRESENCE DANS LES SUJETS
% -----------------------------------------------------------------------

\begin{tcolorbox}[
  ficheSection=colPresence,
  title={Présence dans les sujets analysés (2013-2026)},
]
\begin{center}
\footnotesize\renewcommand{\arraystretch}{1.15}
\begin{tabular}{lccccccccccccc}
\toprule
\textbf{Année} & 13 & 14 & 15 & 16 & 17 & 18 & 19 & 20 & 21 & 22 & 23 & 24 & 25 \\
\midrule
E1 (commune)   & $\checkmark$ & - & - & - & - & - & - & - & - & - & - & - & $\checkmark$ \\
E2 discipl. IE & - & - & - & - & - & - & - & - & - & $\checkmark$ & $\checkmark$ & - & $\checkmark$ \\
\bottomrule
\end{tabular}
\end{center}
\smallskip
\footnotesize\textbf{Sous-thèmes fréquents :}
hacheur abaisseur - courant et ondulation dans $L$ (3/13) - rapport cyclique et tension moyenne (4/13) - diode de roue libre (4/13) - IGBT/MOSFET en convertisseurs (3/13) - diode de protection CAN (2/13)

\medskip
\textbf{Ordres de grandeur (repères rapides).}
\begin{itemize}[leftmargin=1.3em]
  \item Réseau BT FR : tension phase-neutre typique \SI{230}{\volt}, tension phase-phase typique \SI{400}{\volt}, fréquence \SI{50}{\hertz}.
  \item Courants d'étude fréquents en sujet : de quelques \si{\ampere} à quelques dizaines de \si{\ampere}.
  \item En conversion d'énergie, un rendement attendu reste généralement entre 0{,}8 et 0{,}98 selon la technologie.
\end{itemize}

\medskip
\textbf{Synthèse express (à refaire sans support).}
\begin{enumerate}[leftmargin=1.4em]
  \item Poser les hypothèses du modèle et les conventions de signe.
  \item Écrire 2 à 4 relations clés avec unités.
  \item Réaliser un calcul numérique complet et vérifier l'ordre de grandeur.
  \item Justifier le résultat par un argument technique de niveau CAPET.
\end{enumerate}

\end{tcolorbox}

\CAPETENDFILE
